\documentclass[a4paper,12pt]{article}

\usepackage{alltt, fancyvrb, url}
\usepackage{graphicx}
\usepackage[utf8]{inputenc}
\usepackage{float}
\usepackage{hyperref}
\usepackage[italian]{babel}
\usepackage[table]{xcolor}
\usepackage{array}
\usepackage{multirow}
\usepackage{titlesec}
\usepackage{listings}
\usepackage[italian]{cleveref}
\usepackage{acronym}

\sloppy
\setlength{\parindent}{0pt}

\title{\textbf{Project Scoping Meetings}}
\author{}
\date{}

\begin{document}
\maketitle

Partecipanti: Project Manager, team del cliente, membri principali del team di sviluppo.

\section*{Agenda meeting 1}
\begin{itemize}
    \item Introduzione;
    \item scopo del meeting;
    \item descrizione dello stato corrente;
    \item descrizione dell'opportunità di business;
    \item descrizione degli obiettivi che si vogliono raggiungere;
    \item discussione delle Conditions of Satisfaction (CoS);
    \item discussione dei criteri di successo degli obiettivi;
    \item discussione del "gap" tra lo stato corrente e lo stato desiderato;
    \item bozza del POS;
    \item aggiornamento alla prossima riunione per discutere dei requisiti (RBS).
\end{itemize}

\end{document}